\documentclass[11pt]{article}
\usepackage[utf8]{inputenc}
\usepackage{hyperref}
\usepackage{longtable}
\title{AegisScan Project Report}
\author{Ark1414}
\date{\today}

\begin{document}
\maketitle
\begin{abstract}
AegisScan is an agent-less Windows system vulnerability and network scanner designed to identify misconfigurations, insecure services, missing patches, weak firewall rules and unsafe system parameters without installing any software on the target host.
\end{abstract}

\section*{Chapter 1: Introduction}
\subsection*{Overview}
Modern organizations rely heavily on Windows-based infrastructure. AegisScan aims to provide low-overhead, agent-less posture checks.

\subsection*{Objectives and Scope}
Develop an agent-less scanner, detect open services, provide API+UI, and ensure scalability.

\section*{Chapter 2: Background}
\subsection*{Literature Survey}
Existing commercial and open-source scanners provide extensive features but often require agents or subscriptions. Agent-less designs reduce deployment friction.

\section*{Design and Implementation}
Details available in the Markdown report.

\section*{Conclusion}
AegisScan is a modular, lightweight scanner for Windows environments.

\vfill
\small{Generated from project notes.}
\end{document}
